%!TEX root = ../../thesis.tex

\section{Corpus Actions}\label{sec:targets-corpus-actions}

The corpus action model targets are three C parsers, each driven by a persistent
harness parsing a fixed-size seed input in a loop (\cref{sec:corpus-level-actions}).
\emph{cJSON} is run first as an ablation to determine a suitable set of parameters.
\emph{libxml2} and \emph{picohttpparser} then use this configuration for their own campaigns.

\subsection{cJSON}\label{sec:target-cjson}

cJSON 1.7.19 is a single-file recursive-descent JSON parser written in C.
It is the target of our detailed ablation study.
We provide a harness around the \code{cJSON\_ParseWithLength} function, which takes
a buffer and a length and tries to parse the buffer content as JSON.
The harness is shown in~\cref{lst:target-cjson}.
It continuously reads a single 16-byte seed input from standard in and then
tries to parse it.
A single episode then consists of a fixed number of executions of the loop, terminated
by the environment outside the harness.
Afterwards, the state of the emulator is reset and a new episode started.

\begin{ccode}[float=htbp,
  caption={[The cJSON harness]%
    The persistent cJSON harness.
    \code{read\_record} reads a full 16-byte seed and then calls the cJSON parser.},
  label={lst:target-cjson}]
#define INPUT_LENGTH 16

int main(void) {
    char buf[INPUT_LENGTH];

    for (;;) {
        if (read_record(buf) != 0)
            break;

        cJSON *root = cJSON_ParseWithLength(buf, INPUT_LENGTH);
        cJSON_Delete(root);
    }

    return 0;
}
\end{ccode}

In contrast to other targets, we include a detailed ablation study with various parameters.
The different parameter arms are listed in~\cref{tab:cjson-arms}.

\begin{table}[htbp]
  \centering
  \caption[cJSON ablation arms]{%
    The ablation arms of the cJSON study.
    Every arm runs under four different seeds for 200 learner iterations.
  }
  \label{tab:cjson-arms}
  \footnotesize
  \begin{tabularx}{\textwidth}{@{}ll>{\raggedright\arraybackslash}X>{\raggedright\arraybackslash}X@{}}
    \toprule
    Arm & Factor & Baseline & Variation \\
    \midrule
    \code{NOCOV}
      & coverage bitmaps & two 4096-wide maps & zeroed maps \\
    \code{NOSEED}
      & seed bytes & the 16/32 seed bytes & zeroed seed \\
    \code{EDGE}
      & coverage granularity & block hits & edge hits \\
    \code{W0}
      & depth-record weight $\omega$ & 1.0 & 0.0, marginal only \\
      \code{T15}/\code{T25}
      & visit-count temperature $T$ & $1 \to 0.5 \to 0.25$
      & constant 1.5/2.5 \\
    \code{ALPHA004}/\code{ALPHA010}
      & root Dirichlet concentration $\xi$ & 1.0 & 0.04/0.10 \\
      \code{SIMS50}/\code{SIMS100}
      & simulations per search & 200 & 50/100 \\
    \code{REUSE4}
      & training iterations per learner iteration & 16 (reuse of $2\times$)
      & 32 (reuse of $4\times$) \\
    \code{SMALL}
      & hidden width & latent width 128 & latent width 64 \\
    \code{CTRL-RAND}
      & action selection & \gls{mcts}
      & uniform random \\
    \bottomrule
  \end{tabularx}
\end{table}

The corpus holds a single 16-byte seed, an episode is bounded at 64 mutations and therefore 128 MDP
steps, and 128 environments run in parallel (\cref{sec:impl-corpus}).
Every reset rebuilds the corpus from uniformly random bytes, so the agent has to
apply mutations to reach valid JSON syntax on its own.
The reward is the coverage an execution adds to the corpus union plus a bonus per level by
which it out-parses the deepest execution of its episode (\cref{sec:mdp-rewards}).
Because the target executes additional blocks for each correctly guessed
character and therefore rewards guesses in proportion to the length of the
correct prefix, it can successively learn the password character-by-character.

\subsection{libxml2}\label{sec:target-libxml2}

libxml2 v2.15.3. is an XML parser commonly used in Linux.
We disable threads for deterministic execution as well as unreachable and out-of-scope
features like \code{ICONV}, \code{MODULES} \code{CATALOG}, \code{HTML} and \code{XPATH}.
The library entry point the harness makes use of is \code{xmlReadMemory}.
It is structured almost identically as~\cref{lst:target-cjson}.
Instead of 16 bytes, we chose 32 byte seeds to allow richer seed structure.
We also use parse with recovery mode (\code{XML\_PARSE\_RECOVER}).
This causes the parser to ignore errors and instead return a partial tree of successfully parsed input,
yielding a richer coverage signal than simple early abort.

Due to its higher block count, we increase the block hit tracker bucket count
to \num{16384} to reduce collisions.
We do not include a random control baseline.

\subsection{picohttpparser}\label{sec:target-picohttpparser}

picohttpparser is a header-only HTTP/1.x request parser.
We use the latest version from the master branch with commit hash $\texttt{f4d94b4}$.
The setup is identical to libxml2 except the block count of \num{4096}, since this is a smaller target.
